\documentclass[11pt]{article}
\usepackage[a4paper,margin=0.82in]{geometry}
\usepackage{parskip}
\usepackage{ragged2e}
\usepackage[hidelinks]{hyperref}

\pagestyle{empty}
\setlength{\parindent}{0pt}
\setlength{\parskip}{0.48em}

\begin{document}

\begin{flushright}
Zhongli Wang \\
Astastrasse 22, 33617 Bielefeld, Germany \\
+49 160 9541 6015 \quad \href{mailto:nkwangzhongli@gmail.com}{nkwangzhongli@gmail.com} \\
\today
\end{flushright}

\vspace{-0.3cm}

\noindent
Prof. Dr. Ariel Dora Stern \\
Chair for Digital Health, Economics, and Policy \\
Hasso Plattner Institute \\
Prof.-Dr.-Helmert-Str. 2--3 \\
14482 Potsdam, Germany

\vspace{0.2cm}

\noindent
\textbf{Application for Postdoctoral Researcher in Applied Econometrics} \\
\textbf{AI in Clinical Practice}

\vspace{0.15cm}

\noindent Dear Prof. Stern and Members of the Hiring Committee,

I am applying for the Postdoctoral Researcher position in Applied Econometrics: AI in Clinical Practice. I am completing a Ph.D. in Economics at Bielefeld University and Universit\`a Cattolica del Sacro Cuore, with my defense scheduled for September 21, 2026, and I would be available to start in November 2026. My research combines econometric reasoning, computational modeling, machine learning, and the economics of technology adoption.

HPI's Digital Health, Economics, and Policy group particularly appeals to me because it studies new health technologies through econometrics and applied data science. Evaluating an enterprise deployment of clinical decision support is a compelling opportunity to identify not only whether AI changes healthcare delivery, but also how adoption, organizational context, and heterogeneous responses shape its effects. I would value contributing to this work across HPI, Mount Sinai, and MIT while helping coordinate a shared empirical research agenda.

I would bring three complementary strengths to the projects:

\begin{itemize}
    \setlength{\itemsep}{0.05em}
    \setlength{\topsep}{0em}
    \item \textbf{Empirical and computational analysis:} I use Python, R, Stata, simulation-based econometrics, and machine learning to structure data, estimate models, test sensitivity, and communicate results reproducibly.
    \item \textbf{AI and technology research:} My published work on financial fraud detection applies XGBoost to firm-level financial indicators, while a collaborative paper develops agentic research workflows using LLMs and Microsoft's AutoGen.
    \item \textbf{Independent, interdisciplinary delivery:} My doctoral research links industrial organization, heterogeneous-agent modeling, and technology adoption. International coauthorship and conference presentations have trained me to align assumptions, methods, and evidence across disciplines.
\end{itemize}

My empirical work has focused on economic and financial data rather than EHR or claims data, and my econometric specialization to date has been simulation-based. I would therefore enter the position with a clear learning agenda in clinical data structures and design-based causal inference, supported by a strong foundation in statistics, programming, model validation, and reproducible research. I am ready to apply that foundation to matching, difference-in-differences, and experimental evaluations under the team's mentorship.

I would welcome the opportunity to discuss how my background can support rigorous evaluations of clinical AI and the development of follow-on studies, surveys, and randomized trials. Thank you for your consideration.

Sincerely,
\\[0.1cm]
Zhongli Wang

\end{document}
